% ==============================================================================
\chapter{Espacios conexos} \justifying
% ==============================================================================
\section{Espacios conexos y separaciones}

Sea $X$ un espacio topológico.

\begin{defi}
Una \textbf{separación} de $X$ es una pareja $(A, B)$ tal que $A, B \subseteq X$ son abiertos (o cerrados) no vacíos verificando que $X = A \cup B$ y $A \cap B = \emptyset$. Es decir, $A$ y $B$ son complementarios.
\end{defi}

\begin{ejem}
$(\emptyset, X)$ es una separación de $X$, denominada la \textbf{separación trivial}.
\end{ejem}

\begin{defi}
Un espacio topológico $X$ es \textbf{conexo} si no admite separaciones no triviales. Es decir, si $X = A \cup B$ con $(A, B)$ una separación, entonces $A = X$ y $B = \emptyset$ (o viceversa).
\end{defi}

\begin{defi}
Un \textbf{subespacio topológico $Y \subseteq X$} es \textbf{conexo} si lo es como espacio topológico con la topología heredada $(Y, \mathcal{T}_X|_Y)$.
\end{defi}

\begin{obs}
Sean $Y \subseteq Z \subseteq X$ subespacios, $Y \subseteq Z$ es subespacio conexo $\Rightarrow Y \subseteq X$ es subespacio conexo. La noción de conjunto conexo es \textbf{absoluta}, no depende del espacio ambiente (aunque la topología sí depende del ambiente, la conexi\'on se preserva: $\mathcal{T}_Z|_Y = \mathcal{T}_X|_Y$).
\end{obs}

\begin{prop}
Sea $Y \subseteq X$ un subespacio. Si $Y = A \cup B$ con $A \cap B = \emptyset$ siendo $A, B$ abiertos (o cerrados) en $Y$, entonces existen $U, V \subseteq X$ abiertos (o cerrados) en $X$ tales que:
\[
A = U \cap Y, \quad B = V \cap Y \quad (\text{o bien } A = \overline{A}^X \cap Y, \, B = \overline{B}^X \cap Y)
\]
De forma que $Y \subseteq U \cup V$ y $(U \cap V) \cap Y = \emptyset = (U \cap Y) \cap (V \cap Y)$.

Si se verifica que $Y \subseteq U$ o $Y \subseteq V \implies Y \subseteq U \cap Y$ o $Y \subseteq V \cap Y \implies A = Y$ o $B = Y$, haciendo que uno de los dos conjuntos sea el total y el otro el vacío.
\end{prop}

\begin{ejem}[No Conexidad de $\mathbb{Q}$]
En $(\mathbb{R}, \mathcal{T}_d)$, consideremos el subespacio $\mathbb{Q} \subset \mathbb{R}$. Sea $c \in \mathbb{I} = \mathbb{R} \setminus \mathbb{Q}$ un número irracional. Los conjuntos:
\[
((-\infty, c) \cap \mathbb{Q}) = A \neq \emptyset, \quad ((c, \infty) \cap \mathbb{Q}) = B \neq \emptyset
\]
definen una separación no trivial de $\mathbb{Q}$, puesto que $\mathbb{Q} = A \cup B$ y $A \cap B = \emptyset$, concluyendo que \textbf{$\mathbb{Q}$ no es conexo}.
\end{ejem}

\begin{obs}
$\mathbb{R}$ sí es conexo (se probará a continuación). Un espacio conexo puede contener subespacios conexos y no conexos (por ejemplo, $[0, 1] \subset (\mathbb{R}, \mathcal{T}_d)$ es conexo, mientras que $[0, 1] \cup [2, 3] \subset \mathbb{R}$ no es conexo). Idem para espacios no conexos.
\end{obs}

% ==============================================================================
\section{Conexi\'on en $\bb R$}
% ==============================================================================

\begin{prop}
El intervalo $[0, 1] \subset \mathbb{R}$ es conexo (con la topología heredada de la usual $\mathcal{T}_d$).
\end{prop}

\begin{dem}
Para demostrarlo, supongamos que $[0, 1]$ admite una separación no trivial. Es decir, existen $U, V \subseteq \mathbb{R}$ abiertos en $\mathbb{R}$ tales que:
\[
[0, 1] \subseteq U \cup V, \quad U \cap V \cap [0, 1] = \emptyset, \quad U \cap [0, 1] \neq \emptyset, \quad V \cap [0, 1] \neq \emptyset.
\]
Supongamos sin pérdida de generalidad que $0 \in U$ (luego $0 \notin V$). Como $U$ es abierto en $\mathbb{R}$, existe $\varepsilon > 0$ tal que $[0, \varepsilon) \subseteq U$.

Consideremos el conjunto:
\[
T = \{ x \in [0, 1] : [0, x] \subseteq U \} \neq \emptyset.
\]
$T$ es un conjunto de números reales no vacío (pues $0 \in T$) totalmente ordenado y acotado superiormente por $1$. Por el axioma del supremo en $\mathbb{R}$, existe $r = \sup T \in [0, 1]$. Estudiemos la pertenencia de $r$:

\begin{itemize}
    \item Si $r \in U$: Como $U$ es abierto en $\mathbb{R}$, existe $\varepsilon' > 0$ tal que $[r, r + \varepsilon') \subseteq U$. Como $r = \sup T$, existe algún $x \in T$ cercano a $r$, lo que implica $[0, r + \varepsilon') \subseteq U$, de donde $r + \varepsilon' \in T$, contradiciendo que $r = \sup T$.
    \item Si $r \in V$: Como $V$ es abierto en $\mathbb{R}$, existe $\varepsilon'' > 0$ tal que $(r - \varepsilon'', r] \subseteq V$. De nuevo por ser $r = \sup T$, existiría un elemento $x \in T$ tal que $x > r - \varepsilon''$, lo que implicaría que $[0, x] \subseteq U$, resultando en una intersección no vacía de $U$ y $V$ en $[0, 1]$, lo cual es una contradicción.
    \item Si $r=1\Rightarrow r\in U$.
\end{itemize}

Por lo tanto, $[0,1]\subset U$ o $[0,1] \subset V$ la separación no trivial no puede existir, de modo que $[0, 1]$ es conexo.
\end{dem}

\begin{prop}
Todos los intervalos de la forma $[a, b]$, $(a, b)$ y $[a, b)$ son conexos en $\mathbb{R}$.
\end{prop}

\begin{dem}
Se sabe que $[0, 1] \simeq [a, b]$ son homeomorfos mediante una aplicación lineal. Dado que la conexi\'on es una propiedad topológica conservada por homeomorfismos, el intervalo $[a, b]$ es conexo.

Para intervalos generales $Y = [a, b]$, $(a, b)$ o $[a, b)$, si suponemos que $Y \subseteq U \cup V$ es una separación no trivial con $U \cap V \cap Y = \emptyset$, basta tomar un subintervalo cerrado $[c, d] \subset Y$. Como $[c, d]$ es conexo, debe cumplirse $[c, d] \subseteq U$ o $[c, d] \subseteq V$ de manera excluyente, obteniéndose por extensión que $Y$ es conexo.
\end{dem}

\begin{teo}
La recta real $\mathbb{R}$ es conexa.
\end{teo}

\begin{dem}
Se ha demostrado previamente que $(0, 1) \simeq \mathbb{R}$ son homeomorfos. Al ser $(0, 1)$ un intervalo conexo y preservarse la conexi\'on por homeomorfismos, concluimos que $\mathbb{R}$ es un espacio conexo.
\end{dem}

\begin{defi}
Un subconjunto $I \subseteq \mathbb{R}$ es un \textbf{intervalo} si para todo $c, d \in I$ con $c < d$, se satisface que $[c, d] \subseteq I$.
\end{defi}

\begin{prop}
Los únicos subespacios conexos de $\mathbb{R}$ son los intervalos.
\end{prop}

\begin{dem}
Falta demostrar que los intervalos son los \emph{únicos} subespacios conexos de $\mathbb{R}$.

Sea $Y \subseteq \mathbb{R}$ un subespacio conexo. Vamos a probar que $Y$ debe ser un intervalo. Sean $c, d \in Y$ con $c < d$. Queremos ver que $[c, d] \subseteq Y$.

Supongamos por reducción al absurdo que no es un intervalo. Entonces existe algún punto $z \in \mathbb{R}$ tal que $c < z < d$ con $z \notin Y$. 

Consideremos los conjuntos abiertos en $\mathbb{R}$:
\[
U = (-\infty, z), \quad V = (z, \infty)
\]
Notemos que $Y \subseteq U \cup V = (-\infty, z) \cup (z, \infty)$ ya que $z \notin Y$. Además por ser $Y$ conexo
\[
c \in U \cap Y \implies U \cap Y \neq \emptyset
\]
\[
d \in V \cap Y \implies V \cap Y \neq \emptyset
\]
\[
(U \cap Y) \cap (V \cap Y) = (U \cap V) \cap Y = \emptyset \cap Y = \emptyset
\]
Por lo tanto, $(U \cap Y, V \cap Y)$ constituye una separación no trivial del subespacio $Y$, lo cual contradice la hipótesis de que $Y$ es conexo.

Concluimos que $Y$ debe ser un intervalo.
\end{dem}
\section{Propiedades}
\begin{prop}
    Sea $Y \subset X$ subespacio topol\'ogico conexo y $(A,B)$ separaci\'on de $X$. Entonces $Y \subset A$ \'o $Y \subset B$. Es decir, que si $Y \subset X = A\sqcup B$ ($\sqcup$ como conjuntos), entonces $Y \subset A$ \'o $Y\subset B$
\end{prop}
\begin{dem}
    $Y = (Y\cap A ) \cup (Y \cap B)$ es una separaci\'on $Y$, que es conexo, y $Y\cap A $ es abierto y $Y\cap B$ es abierto, entonces $Y\cap A = Y$ \'o $Y\cap B = Y$, entonces $Y \subset A$ \'o $Y\subset B$.  
\end{dem}
\begin{prop}
    Si $A,B \subset X$ subconjunto con $A \cup B = X$ y $\overline{A}\cap B = \overline{B}\cap A = \emptyset$ entonces $(A,B)$ es separaci\'on de $X$.
\end{prop}
\begin{dem}
    Como $A\subset \overline{A}$ autom\'aticamente $A \cap B = \emptyset$. Veamos que son cerrados. Como $\overline{A} = \bar A \cap X = \bar A \cap (A\cup B) = (\bar A \cap A) \cup (\bar A \cap B) = A \cup \emptyset = A$. De la misma forma con $B$, y como ambos son cerrados, es una separaci\'on.
\end{dem}
\begin{cor}
Si $X$ es conexo, entonces no existen subconjuntos $A,B \subset X$ con $A \cup B = X$ y $\bar A \cap B = \bar B \cap A = \emptyset$.
\end{cor}

Sea $X$ espacio topol\'ogico y $\{X_i \subset X\}_{i \in I}$ subespacios conexos.
\begin{ejer}
    Se verifica que $\bigcup X_i$ es conexo si $\bigcap X_i \neq \emptyset$. $\bigcup X_i$ es conexo si para todo $i, j\in I$ $X_i \cap X_j \neq \emptyset$. La segunda implica la primera. 
\end{ejer}
Demostraremos la primera
\begin{dem}
    $\bigcup X_i \subset A \cup B$ con $A,B$ abiertos, $A\cap B \cap (\bigcup X_i) = \emptyset$ lo que da lugar a una separaci\'on. Para cada $j\in I$, $X_j \subset \bigcup X_i \bigcup A \cup B$ y como es conexo y tiene separaci\'on, $X_j \subset A$ \'o $X_j \subset B$. 

    Sea $x_o \in \bigcap X_i \neq \emptyset \Rightarrow x_o \in X_i\subset A \cup B \, \forall i$ por tanto $x_o \in A$ \'o $x_o \in B$. Supongamos que $x_o \in A \Rightarrow x_o \nin B$, entonces si para alg\'un $j\in I$, $x_o \in X_j \in B$, como $x_o \in X_i \forall i$, $x_o \in B$, contradicci\'on con $x_o \nin B$, luego $X_j \subset A \forall \, j\in I$ por lo que $\bigcup X_j \subset A$. 
\end{dem}
\begin{prop}
    Sea $Y \subset X$ subespacio conexo. Para todo $Z\subset X$ con $Y\subset Z \subset \overline{Y}$ se verifica que $Z$ es conexo. En particular, si $Y$ es conexo $\Rightarrow \overline{Y}$ es conexo.
\end{prop}
\begin{dem}
    Sea $(A,B)$ una separaci\'on de $Z$, $Z = A\cup B$ y $A\cap B=\emptyset$. Supongamos $A$ y $B$ cerrados de $Z$, y $A = \overline{A} \cap Z$ y $B = \overline{B}\cap Z$. 
    
    $Y \subset Z = A\cup B$ y como $Y$ es conexo por la primera proposici\'on, $Y \subset A$ \'o $Y \subset B\Rightarrow \overline{Y} \subset \overline{A}$ \'o $\overline{Y} \subset \overline{B}$. Entonces al cortar con $Z$, $\overline{Y}\cap Z \subset \overline{A} \cap Z$ que es $Z \subset A \Rightarrow Z=A$ \'o de la misma forma $Z=B$.  
\end{dem}
\begin{prop}
    Sean $X,Y$ espacios topol\'ogicos, $f: X \longrightarrow Y$ aplicaci\'on continua. Si $A\subset X$ es subespacio conexo, entonces $f(A) \subset Y$ es un subespacio conexo. 

    En particular, si $X$ es conexo, entonces $\Im f$ es conexo.
\end{prop}
\begin{dem}
    Sea $V_1, V_2 \subset Y$ abiertos con $f(A) \subset V_1 \cup V_2$ y $V_1 \cap V_2 \cap f(A) = \emptyset$. 
    $$f^{-1}(f(A)) \subset f^{-1}(V_1 \cup V_2) = f^{-1}(V_1) \cup f^{-1}(V_2)$$
    donde $f^{-1}(V_i)$ son abiertos por ser $f$ continua y adem\'as disjuntos con $f(A)$ por $V_1 \cap V_2 \cap f(A) =\emptyset \Rightarrow f^{-1}(V_1) \cap  f^{-1}(V_2)\cap f^{-1}(f(A)) = \emptyset$ y como $A$ es conexo, 
    $$A \subset f^{-1}(V_1) \text{ \'o } A \subset f^{-1}(V_2) $$
    de forma que $f(A)\subset f( f^{-1}(V_1) )\subset V_1$ \'o para $V_2$. Es decir, $f(A)$ es conexo. 
\end{dem}
\begin{cor}
    [Teorema del valor intermedio] Sea $X$ un espacio topol\'ogico conexo, $f:X \longrightarrow \bb R$ funci\'on continua. Si $a,b \in \Im f$, para todo $c\in \bb R$ con $a\leq c \leq b$ ($c\in [a,b]$), existe $x\in X$ tal que $f(x) = c$ ($c\in \Im f$). 
\end{cor}
\begin{dem}
    $f(X) \subset \bb R$ es conexo, y como los conexos de $\bb R$ son los intervalos, $f(X)=I$ intervalo de $\bb R$, $a, b\in I$ entonces $a,b \in \Im f$ entonces si $c\in [a,b] \subset \Im f$. 
\end{dem}
\section{Producto directo y uni\'on disjunta}
\begin{obs}
    Un subespacio conexo si y solo si los \'unicos abiertos-cerrados que tiene son $X, \emptyset$. Si tenemos un abierto-cerrado no trivial tenemos una separaci\'on pues su complementario tambi\'en es abierto-cerrado. 
\end{obs}
$\{X_i\}$ espacios topol\'ogicos, se tiene el espacio topol\'ogico $\bigsqcup X_i$, solo falta ver que $X_i$ son abiertos o cerrados. Lo son porque $\bigsqcup X_i$ tiene la topolog\'ia final de las $\varphi_j$, de forma que $X_j\subset \bigsqcup X_i$ son abiertos y cerrados. Luego $\bigsqcup X_i$ no es conexo aunque $X_j$ lo sean. 